\documentclass[10pt,a4paper]{amsart}
\usepackage[margin=19mm]{geometry}
\usepackage{amsmath,amssymb,mathtools}
\usepackage[scr=rsfs,cal=euler]{mathalfa}
\usepackage{xfrac}
\usepackage[unicode,hidelinks,bookmarks=false]{hyperref}
\newcommand{\ie}{i.e.\ }
\newcommand{\cf}{cf.\ }
\newcommand{\resp}{respectively\ }
\newcommand{\Subsection}{\subsection}
\providecommand{\nobreakdash}{\nobreak\hbox{-}}
\setlength{\emergencystretch}{2em}
\multlinegap0pt
\usepackage{amssymb}
\usepackage{mathtools}
\usepackage[shortlabels]{enumitem}
\usepackage{xcolor}
\newtheorem{proposition}{Proposition}
\newcommand{\Gal}{\operatorname{Gal}}
\newcommand{\Q}{\mathbb Q}
\newcommand{\Qp}{\mathbb Q_p}
\newcommand{\Z}{\mathbb Z}
\newcommand{\Zp}{\mathbb Z_p}
\title{Galois representations ramified at one prime via relative deformation theory}
\author{Anwesh Ray}
\date{Source: arXiv:2609.03954v1 (CC BY 4.0)}
\begin{document}
\maketitle

\noindent\emph{Source and licence.} Galois representations ramified at one prime via relative deformation theory. Version arXiv:2609.03954v1 (CC BY 4.0), distributed by arXiv under the Creative Commons Attribution 4.0 International licence, http://creativecommons.org/licenses/by/4.0/, as recorded in the arXiv licence metadata for this exact version and shown on the abstract page https://arxiv.org/abs/2609.03954v1. Full licence notice: \texttt{licenses/arxiv-2609.03954-CC-BY-4.0.txt}.\par\smallskip

\noindent\textit{Editorial scope.} Counted unit: the article's proposition prop:cyclotomic-specialization (source0.tex lines 825--831). The article's complete original proof of it is delivered with it (source0.tex lines 833--970, one \texttt{proof} environment of 138 lines). No mathematics is written, repaired or re-derived: the statement and the proof are the article's own lines, in the article's own notation, and no retained line is substituted or re-flowed.  Retained uncounted as the source-local context the proof needs: lines 808--810.  Nothing was omitted from any retained span, so the package carries no omission record.  No retained line contains a citation, so the wrapper carries no bibliography and no bibliography entry.  No retained line contains a figure environment, an image-inclusion command or drawing code, so no image is delivered and the package declares no asset dependency.  Editorial formatting only, in addition: an AMS \texttt{amsart} preamble that restates the article's own declarations and macros used by the retained lines, namely the environments \texttt{proposition} on separate counters, together with the standard packages those lines need.  The retained environments print in this document as Proposition 1 in order of appearance, whereas the article prints the same items with the numbering of its own full document.  The complete native source of the article as distributed in the arXiv e-print for this exact version is bundled unchanged as \texttt{sources/209329/source0.tex}.
\begin{equation}\label{eq:weak-leopoldt-H}
 H^2(H,\Qp/\Zp)=0,
\end{equation}

\begin{proposition}\label{prop:cyclotomic-specialization}
For every $m\in\Z$ there is a natural isomorphism
\begin{equation}\label{eq:H2-specialization}
 H^2\!\left(G_{\{p\}},(\Qp/\Zp)(m)\right)^\vee
 \simeq e_mX[f_m].
\end{equation}
\end{proposition}

\begin{proof}
Put $A:=\Qp/\Zp$ and $A_m:=A(m)$. Recall that
$H=\Gal(\Q_{\{p\}}/F_\infty)$ and
$\Gamma=\Gal(F_\infty/\Q)$. Since $F_\infty$ contains all
$p$-power roots of unity, the cyclotomic character is trivial on
$H$. Thus $H$ acts trivially on the Tate-twist factor of $A_m$.
Consequently, for every $q\geq0$, there is an identification
$H^q(H,A_m)=H^q(H,A)(m)$ as $\Gamma$-modules. 

We first apply the Hochschild--Serre spectral sequence to the exact
sequence
\[
 1\longrightarrow H\longrightarrow G_{\{p\}}
 \longrightarrow\Gamma\longrightarrow1.
\]
It has second page
\[
 E_2^{a,b}
 =
 H^a\bigl(\Gamma,H^b(H,A_m)\bigr)
 \Longrightarrow
 H^{a+b}(G_{\{p\}},A_m).
\]
We are interested in total degree $2$. There are therefore three
potential terms, namely $E_2^{0,2}$, $E_2^{1,1}$ and $E_2^{2,0}$. The first of these vanishes. Indeed, weak Leopoldt \eqref{eq:weak-leopoldt-H} for the
cyclotomic extension gives $H^2(H,A)=0$. Since the Tate twist does
not change the action of $H$, it follows that $H^2(H,A_m)=0$, and
hence $E_2^{0,2}=0$.

The third term also vanishes. Recall that
$\Gamma=\Delta_p\times\Gamma_1$, where $\Gamma_1\simeq\Zp$ and
$|\Delta_p|=p-1$. The group $\Gamma_1$ has $p$-cohomological
dimension one, while a finite group of order prime to $p$ has
$p$-cohomological dimension zero. Hence
$\operatorname{cd}_p(\Gamma)=1$. Since $A_m$ is a discrete
$p$-primary module, this gives
$E_2^{2,0}=H^2(\Gamma,A_m)=0$.

Thus the only possibly nonzero term in total degree $2$ is
\[
 E_2^{1,1}
 =
 H^1\bigl(\Gamma,H^1(H,A_m)\bigr).
\]
Since $H^1(H,A_m)=Y(m)$, this term is
$H^1(\Gamma,Y(m))$. Moreover, no differential can enter
$E_2^{1,1}$, since its possible source would have negative first
index, and no differential can leave it: the first possible target
is $E_2^{3,0}$, which vanishes because
$\operatorname{cd}_p(\Gamma)=1$. Therefore
$E_\infty^{1,1}=E_2^{1,1}$. As all the other graded pieces of the
Hochschild--Serre filtration on $H^2(G_{\{p\}},A_m)$ vanish, we
obtain a canonical isomorphism
\begin{equation}\label{eq:HS-H2}
 H^2(G_{\{p\}},A_m)
 \simeq H^1(\Gamma,Y(m)).
\end{equation}

We next compute the group on the right. Apply Hochschild--Serre to
the exact sequence $1\to\Gamma_1\to\Gamma\to\Delta_p\to1$. Since
$|\Delta_p|=p-1$ is prime to $p$, one has
$H^a(\Delta_p,D)=0$ for every $a>0$ and every discrete
$p$-primary $\Delta_p$-module $D$. Indeed, multiplication by
$|\Delta_p|$ is an automorphism of $D$, so the usual averaging
operator makes the invariants functor exact. It follows that
\[
 H^1(\Gamma,Y(m))
 \simeq H^1(\Gamma_1,Y(m))^{\Delta_p}.
\]

Let $\gamma$ be the fixed topological generator of $\Gamma_1$ and
put $u=\chi(\gamma)$. For any discrete $p$-primary
$\Gamma_1$-module $D$, continuous cohomology of
$\Gamma_1\simeq\Zp$ gives
$H^1(\Gamma_1,D)=D/(\gamma-1)D$. On the twisted module $Y(m)$,
the element $\gamma$ acts as $u^m\gamma$, where the second
$\gamma$ denotes its original action on $Y$. Consequently,
\[
 H^1(\Gamma_1,Y(m))
 \simeq Y/(u^m\gamma-1)Y.
\]
For brevity, denote this quotient by $Q_m$. Thus
$H^1(\Gamma,Y(m))=Q_m^{\Delta_p}$.

We now dualize. Since $X=Y^\vee$, the Pontryagin dual of $Q_m$ is
the annihilator in $X$ of $(u^m\gamma-1)Y$. Under the usual
contragredient action on the Pontryagin dual, the adjoint of
$\gamma$ acting on $Y$ is $\gamma^{-1}$ acting on $X$. Therefore
\[
 Q_m^\vee
 =
 \ker\bigl(u^m\gamma^{-1}-1:X\longrightarrow X\bigr).
\]
Multiplying the operator by the unit $u^{-m}$ does not change its
kernel. Hence
\[
 Q_m^\vee
 =
 \ker(\gamma^{-1}-u^{-m}:X\longrightarrow X)
 =
 X[f_m],
\]
where, by definition, $f_m=\gamma^{-1}-u^{-m}$.

It remains to keep track of the action of the finite group
$\Delta_p$. This is the point at which the idempotent $e_m$ enters.
For $\delta\in\Delta_p$, the action on the twisted module $Y(m)$ is
the original action on $Y$ multiplied by
$\chi(\delta)^m=\omega(\delta)^m$. Therefore an element of $Q_m$
is fixed by $\Delta_p$ after twisting precisely when, with respect
to the original action on $Y$, it belongs to the
$\omega^{-m}$-eigenspace. In terms of the standard idempotents,
this says
\[
 Q_m^{\Delta_p}=e_{-m}Q_m.
\]

Pontryagin duality reverses the finite character. More explicitly,
if $y$ transforms under $\Delta_p$ by $\omega^{-m}$, then the
contragredient action on a functional pairing nontrivially with
$y$ is by $\omega^m$. Equivalently, under the involution
$\delta\mapsto\delta^{-1}$ on $\Zp[\Delta_p]$, the idempotent
$e_{-m}$ is carried to $e_m$. It follows that
\[
 (Q_m^{\Delta_p})^\vee
 =
 (e_{-m}Q_m)^\vee
 \simeq e_mQ_m^\vee.
\]
Using $Q_m^\vee=X[f_m]$ and the fact that $f_m$ commutes with the
$\Delta_p$-action, we finally obtain
\[
 H^1(\Gamma,Y(m))^\vee
 \simeq e_mX[f_m].
\]
Combining this with \eqref{eq:HS-H2} proves
\eqref{eq:H2-specialization}.
\end{proof}

\end{document}
